\cotefile{85}
\scanpage{44}

Sous-groupes finis de \deleted{\unsure{Aut}} $\underline{\mathrm{Aut}}_{\!X}\deleted{(X)}=\mathbb A$ ($X$ droite projective sur corps \unsure{alg} \inserted{de car.~$0$})

\unsure{1.)|1)} Sous-groupes $\Gamma^{+},\Gamma\subset\mathbb A$ qui sont des formes de $\unsure{A_4|\mathcal A_4|\mathfrak A_4}$ ou $\unsure{\mathfrak S_4|\mathcal S_4}$. Se correspondent biunivoquement

\[
\Gamma^{+}\longmapsto\Gamma=\mathrm{Norm}_{\unsure{G|A|\mathbb A}}(\Gamma^{+})\qquad
\Gamma^{+}=[\Gamma,\Gamma]
\]

\uline{NB} $\mathrm{Aut}(A_4)\xleftarrow{\sim}\mathrm{Aut}(\mathfrak S_4)\xleftarrow{\sim}\mathfrak S_4$ On aura : la cat. des ensembles à $4$ él., celle des formes de $A_4$, celle des formes de $\mathfrak S_4$, sont équivalentes

\diagram{\begin{tikzcd}[ampersand replacement=\&,column sep=small]
 F(\mathfrak S_4)\arrow[rr,dash,"\sim"]\arrow[dr,dash,"\sim"'] \& \& F(A_4)\arrow[dl,dash,"\sim"]\ {\deleted{\unsure{(Ens)_4}}} \\
 \& (\mathrm{Ens})_4 \&
\end{tikzcd}}

\[
\begin{aligned}
E\in(\mathrm{Ens})_4,\ &\longmapsto\Gamma=\mathfrak S_E,\ \Gamma^{+}=\mathfrak S_E^{+}\\
\Gamma^{+}\in F(A_4)&\longmapsto\Gamma=\mathrm{Aut}(\Gamma^{+}),\ \text{$E\unsure{=}$ ens des ss-gpes d'indice $4$}\\
&\hspace{6em}\text{(d'ordre $3$) de $\Gamma^{+}$}\\
\Gamma\in F(\mathfrak S_4)&\longmapsto\Gamma^{+}=[\Gamma,\Gamma],\quad E=\text{ens des ss-gpes $\inserted{\gamma}$ d'ordre $6$ (indice $4$)}\\
&\hspace{3em}\text{de $\Gamma$}\simeq\text{ ens des ss-gpes}\\
&\hspace{4em}\gamma^{+}\text{ d'ordre $3$ de $\Gamma$}\\
&\quad[\,\gamma^{+}=[\gamma,\gamma],\ \gamma=\mathrm{Norm}_{\Gamma}(\gamma^{+})\,]
\end{aligned}
\]

Donc sur schéma $S$

\[
\mathrm{Rev}_4(S)\simeq F(S,\mathfrak S_4)\simeq F(S,A_4)
\simeq\mathrm{Tors}(\mathfrak S_{4S})\simeq\text{\uline{Cubes orientés sur $S$}}
\]

\deleted{\unsure{$\simeq$}} \uline{Re NB} En car $0$ (\unsure{inutile} ?) Si $\Gamma^{+}$ ou $\Gamma\subset\deleted{X,\ \text{alors}}\inserted{G=\underline{\mathrm{Aut}}_{\!X}},$ alors [tout autom de $\mathbb A$ est int, $\mathbb A\mathrel{\overset{\mathrm{int}}{\simeq}}\mathrm{Aut}(\mathbb A),$ et] tout autom de $(\mathbb A,\Gamma^{+})$ ou $(\mathbb A,\Gamma)$ qui normalise $\Gamma^{+}$ ou $\Gamma$ provient de $\Gamma$ : $\Gamma\xrightarrow{\sim}\underline{\mathrm{Aut}}_{(X,\Gamma)},\ \Gamma\simeq\underline{\mathrm{Aut}}_{(X,\Gamma^{+})}$

Donc la cat. des couples $(\Gamma,X)$ $(\Gamma^{+},X)$ est équivalente à celle des torseurs sous $\mathfrak S_{4S}$, i.e. aux cat. précédentes -- donc aussi à celle des \uline{cubes orientés} sur $S$, ou celle des \uline{tétraèdres} (pas \deleted{\unsure{néce.}} orientés) sur $S$. Dans ce

% archive sheet number: 43
